\section{Discussion and limitations}\label{sec:discussion}
The four studies suggest an order of priorities for a multi-UAV replanner under uncertain execution. The planning
model matters most. Moving every deadline by a margin that grows with the lead time raised the utility by more than any
choice of stopping rule or trigger, on three independent sets of held-out seeds and also without search. The trigger
comes second. Periodic replanning helped every method on the nominal model, but only by \val{KDabsAblTrig}\,pp on the
chance-constrained model, which suggests that both correct the same accumulated deviations, one before and one after
they occur. The stopping rule comes last. On the proposed model, it was equivalent to a matched fixed budget in
study~2 and added \val{XDabsFullCC}\,pp in study~4, and its main value lies in avoiding the losses of long budgets on
slow planners. The metareasoning rules behaved in the same way, probably because the improvement curve of this ALNS
flattens within a few tens of milliseconds.

Study~4 separates the shape of the margin from the way it is computed. The constant buffer was clearly worse than every
margin that grows with the lead time, whereas the calibrated effective deadline, a proportional buffer with a tuned
slope and moment propagation along the route were within \val{XRangeScaled}\,pp of each other. This agrees with the
error structure of Section~\ref{sec:cc}. The spread grows almost linearly with the lead time, so a margin $c\,h$ with a
suitable slope captures most of it, and the intercept of the calibrated model adds little. The calibrated model does
not improve on the tuned buffer. Its value lies in setting the slope from flight data and in a predicted probability of
meeting the deadline that ranked the tasks by their risk (Supplementary Section~S7). Under high disturbances, the slope that the calibrated model
implied was almost the one tuned on utility. Moment propagation evaluates the actual legs of a route and was slightly
better by an amount within the noise, but it reduced the search rate. A planner with a tighter budget or larger fleets
would feel this cost more, whereas a planner with a generous budget may prefer the more detailed model.

The fallback was the largest single effect of studies~3 and~4. A hard chance constraint discards every task whose
probability of being served on time falls below $\Phi(z)$, even when a vehicle has nothing better to do. Study~4 shows
that this holds however the constraint is evaluated. Without fallback, moment propagation served the fewest tasks late
of all models but the least reward, below the nominal model. Keeping the given-up tasks with their nominal deadline
costs punctuality, because many fallback tasks are served late. A soft formulation that weights every task by its
probability of being served on time would make this trade explicit and could be tuned to the cost of a late arrival in a
given application.

%S4C-DISC

Several limitations bound these conclusions. The error model is linear in the lead time and fitted per disturbance
level, so the level is assumed to be known, and $z$ was selected per level on development missions. Online calibration
within a mission removes this assumption but served \val{XDabsOnlCC}\,pp less reward, and its non-inferiority was not
shown. The chance constraint is imposed on every task separately rather than jointly on a route. The fallback and the
models of study~4 were designed after the held-out results of earlier studies and were therefore tested on new seeds.
Tier C uses three UAVs and fifteen tasks per mission. Software-in-the-loop simulation reproduces the autopilot and the
vehicle dynamics but not the sensors, the radio link or the weather of a field test, and the wind is constant during a
mission. The latency factor emulates a slower planning computer instead of measuring one. Because the stopping rules
read a wall clock, repeated runs are not bit-identical, and the mean utility of a method varied by up to
\val{JRepeatDiffMax}\,pp between two runs of the same missions. Finally, all methods share one ALNS and one family of
scenarios, and other search methods or mission types may change the size of the gains.
